% Research questions — pivoted framing (the DSL kernel *evaluation* problem).
% Canonical framing: ../PIVOT_FRAMING.md. Pre-pivot RQ definitions are in git history.

We organize the study around three research questions.
The first asks whether existing benchmarks can tell an efficient DSL kernel from a
performance-poor one; the second characterizes and explains the performance gap those
benchmarks fail to surface; the third asks how to guide kernel development when a
comprehensive benchmark is infeasible.

\noindent\textbf{RQ1 (Evaluation gap).} Do existing DSL and LLM-generated-kernel
benchmarks distinguish efficient kernels from performance-poor ones, and along which
dimensions---DSL coverage, data type, input shape, hardware, and reward-hacking
prevention---do they fall short?

\medskip\noindent\textbf{RQ2 (Hidden gap and its causes).} For kernels that pass
existing correctness-style benchmarks, how large and structured is the performance gap
to vendor libraries, and which authoring, code-generation, and library-maturity factors
cause it?

\medskip\noindent\textbf{RQ3 (Guidance without a comprehensive benchmark).} Absent a
comprehensive benchmark, can lightweight evaluation heuristics---comparability with the
vendor baseline plus a roofline anchor---together with recurring optimization patterns
reliably flag poor kernels and guide fixes? Is there a single dominant fix, and where
are DSLs genuinely limited?

\medskip\noindent
Our primary measurements span two datacenter architectures, the NVIDIA A100-SXM4-40GB (Ampere, \texttt{sm\_80}) and the NVIDIA GH200 (Grace Hopper, \texttt{sm\_90}), spanning two GPU families.
